\section{Diskussion}
\label{sec:Diskussion}

In dem vorliegenden Experiment wurden akustische Resonatoren erfolgreich als klassische Analogiesysteme für quantenmechanische Zustände untersucht. Die Auswertung zeigt, dass trotz fundamental unterschiedlicher physikalischer Natur der zugrundeliegenden Phänomene – Schallwellen in Gasen versus Wahrscheinlichkeitsamplituden von Elementarteilchen – tiefe mathematische und strukturelle Gemeinsamkeiten existieren.

\subsection{Bestimmung der Schallgeschwindigkeit als Vorbereitung}
Der erste Teil des Experiments diente der Charakterisierung der zylindrischen Resonatoren und der Bestimmung der Schallgeschwindigkeit $c$. Die experimentell ermittelten Werte betragen:
\begin{align}
    c_{50\text{mm}} &= (327{,}10 \pm 22{,}75)\,\text{m/s} \\
    c_{75\text{mm}} &= (329{,}10 \pm 24{,}87)\,\text{m/s}
\end{align}
Diese Werte liegen extrem nah beieinander und stimmen innerhalb ihrer statistischen Unsicherheiten überein. Ein Vergleich mit dem Literaturwert der Schallgeschwindigkeit in trockener Luft bei Raumtemperatur ($20^\circ\text{C}$), welcher bei $c_{\text{Lit}} \approx 343\,\text{m/s}$ liegt, ergibt folgende prozentuale Abweichungen:
\begin{itemize}
     \item $\Delta_{\text{, 50mm}} = \frac{|327{,}10 - 343|}{343} \approx 4{,}64\,\%$
     \item $\Delta_{\text{, 75mm}} = \frac{|329{,}10 - 343|}{343} \approx 4{,}05\,\%$
\end{itemize}
Diese geringfügigen systematischen Abweichungen nach unten lassen sich durch thermische Abweichungen und Fehlstellen an den Zusammensetzungspunkten der Zylinder erklären. Da beide Werte jedoch um weniger als fünf Prozent vom Literaturwert abweichen, bestätigt dieses Ergebnis das theoretische Modell der stehenden Wellen mit starren Randbedingungen (Neumann-Randbedingungen). 

\subsection{Die Wasserstoffanalogie im Kugelresonator}
Im Hauptteil des Versuchs konnte die Winkelabhängigkeit der Druckwellen im Kugelresonator präzise vermessen werden. Die erhaltenen Polarplots für die Eigenfrequenzen bei $2{,}3\,\text{kHz}$, $3{,}7\,\text{kHz}$, $5{,}0\,\text{kHz}$ und $7{,}4\,\text{kHz}$ zeigen eine ziemlich gute Übereinstimmung mit $Y_{l,m}(\theta, \varphi)$. Das zeigt, dass die geometrische Symmetrie eines Hohlraums die mathematische Form der Wellenfunktion vorgibt, unabhängig davon, ob es sich um Schallwellen oder die Aufenthaltswahrscheinlichkeit eines Elektrons im Coulomb-Potential im Atom handelt. 

Es muss jedoch angemerkt werden, dass der radiale Anteil der Helmholtz-Gleichung (Besselfunktionen) fundamental vom radialen Anteil der Schrödinger-Gleichung für das Wasserstoffatom (Laguerre-Polynome) abweicht. Daher lassen sich aus den absoluten Peakpositionen im Übersichtsspektrum keine direkten Rückschlüsse auf die quantenmechanischen Energieeigenwerte (die Hauptquantenzahl $n$) ziehen. Das akustische System ist eine Winkel-Analogie, kein exaktes physikalisches Abbild des Wasserstoffatoms.

\subsection{Aufhebung der Entartung durch Symmetriebrechung}
Ein zentraler Erfolg des Experiments ist der Nachweis der Entartungsaufspaltung. Ohne Modifikationen waren die verschiedenen Drehimpulszustände aufgrund der idealen Kugelsymmetrie bei derselben Resonanzfrequenz entartet, wobei die Zustände mit höherem Drehimpuls $l$ in den Messungen dominierten. Durch das Einbringen der Zwischenringe mit Dicken von $3\,\text{mm}$ bis $9\,\text{mm}$ wurde diese Symmetrie gebrochen. 

Die Messergebnisse in Abbildung 9 zeigen deutlich isolierte Polarplots für die Zustände $l=0$ (isotrope Kugelform) und $l=1$ (Dipolcharakter). Dies stellt eine Analogie zu quantenmechanischen Effekten wie dem Zeeman-Effekt dar, bei denen äußere Felder die Richtungsquantisierung aufheben und Energieniveaus aufspalten. Dass die Entartung für $\pm m$ aufgrund der verbleibenden Rotationssymmetrie um die vertikale Achse theoretisch erhalten bleibt, deckt sich ebenfalls mit den physikalischen Erwartungen.

\noindent Kleinere Abweichungen in den Polarplots (insbesondere Asymmetrien in den Amplitudenmaxima) lassen sich auf experimentelle Unzulänglichkeiten zurückführen. Dazu zählen minimale Justagefehler beim manuellen Einstellen der Winkel in $5^\circ$-Schritten, akustische Streuungen an den Kanten der eingesetzten Ringe sowie Hintergrundgeräusche oder Vibrationen (beispielsweise erkennbar, wenn während der Messung auf den Tisch geklopft wird).